\section{Complexity and Practical Considerations}
\label{sec:complexity}
The measurement pipeline adds no cost to inference: fidelity is computed offline,
once, by comparing baseline and compressed attributions. Computing an attribution
costs one backward pass (saliency, Grad-CAM), $n$ backward passes (IG, here
$n=20$), or $O(\text{patches})$ forward passes (occlusion). Table~\ref{tab:cost}
reports the measured per-image attribution latency on CPU: Grad-CAM and saliency
cost under $10$~ms, occlusion $39$~ms, and IG $134$~ms (dominated by its $n$
backward passes), while a full post-training quantization or pruning pass over
the network takes only a few milliseconds. The corresponding weight-memory
compression ratio is $32/b$, i.e.\ $8\times$ at $4$~bits and $10.7\times$ at
$3$~bits over the FP32 model. EAQ adds only a handful of calibration backward
passes plus $O(L)$ arithmetic for the closed-form allocation of
Theorem~\ref{thm:waterfill}, and---unlike quantization-aware training---requires
\emph{no} retraining, making it suitable for on-device post-training deployment.
All experiments completed on a single commodity CPU with $<4$\,GB of memory,
deliberately matching the constrained hardware that motivates compression; the
same code runs unmodified on the Apple~M3~Pro MPS backend and on CUDA GPUs.

\begin{table}[!t]
\caption{Measured computational cost on CPU. Left: per-image attribution and
per-model compression latency. Right: weight-memory compression ratio vs.\ FP32.}
\label{tab:cost}
\centering
\begin{tabular}{l r @{\hskip 1.2em} c r}
\hline
\textbf{Operation} & \textbf{ms} & \textbf{Bits} & \textbf{Ratio}\\
\hline
Saliency (per img)   & 9.2   & 8 & $4.0\times$\\
Grad-CAM (per img)   & 8.6   & 6 & $5.3\times$\\
Occlusion (per img)  & 39.0  & 5 & $6.4\times$\\
IG, $n{=}20$ (per img) & 134.1 & 4 & $8.0\times$\\
Quantize (per model) & 4.4   & 3 & $10.7\times$\\
Prune (per model)    & 5.4   & 2 & $16.0\times$\\
\hline
\end{tabular}
\end{table}

\section{Discussion, Limitations, and Future Work}
\label{sec:discussion}

\subsection{Interpretation}
Our results reframe a common but untested assumption. Practitioners routinely
validate a model's interpretability at full precision and then deploy a
compressed variant, implicitly assuming the explanations carry over. We show this
assumption is \emph{safe in the graceful regime} ($b\ge5$, $s\le50\%$) and
\emph{unsafe in the collapse regime}, and---critically---that accuracy gives no
warning of the transition. The fidelity--bit-width law
(Theorem~\ref{thm:law}) turns this into an actionable rule: given a tolerable
fidelity loss $\varepsilon$ and an estimate of $\kappa$, deploy at
$b\ge\tfrac12\log_2(\kappa/\varepsilon)$.

The fidelity knee near $b^\star\!\approx\!4.7$ also admits an
information-capacity reading consistent with a \emph{fidelity--plasticity}
trade-off observed in low-bit parameter-efficient fine-tuning
\cite{guo2024lqlora}: once weight precision drops below roughly four bits the
frozen representation is too degraded to support coherent feature extraction, so
downstream capacity (there, adapter rank; here, the informativeness of an
attribution) yields diminishing returns. Below this precision the collapse of
explanation maps into fragmented, off-object noise (Fig.~\ref{fig:qual}) is the
attribution-space signature of the same underlying information bottleneck.

\subsection{Limitations}
Several caveats bound the scope of our claims, and each maps to a concrete
extension of the released protocol (Section~\ref{sec:setup-extended}).
\emph{(i)~Dataset and architecture scope.} The controlled SynthShapes-32
benchmark buys ground-truth masks at the cost of natural-image complexity. However, we have extended this protocol to standard natural-image benchmarks (CIFAR-10 and ImageNette) with ResNet backbones (Section~\ref{sec:setup-extended}), confirming that the graceful-degradation profile, the robustness ordering ($\text{Grad-CAM}\succeq\text{Occlusion}\succeq\text{IG}\succeq\text{saliency}$), and the EAQ advantages hold at that scale as well. Evaluating additional post-training and quantization-aware compressors (e.g.\
GPTQ-, AWQ- and SparseGPT-style schemes), comparing EAQ against established
mixed-precision allocators such as HAWQ \cite{dong2019hawq} and AdaRound
\cite{nagel2020adaround} rather than against uniform
quantization alone, and adding a calibration-size sensitivity study (how few
images $\hat\omega_l$ can be estimated from) together with further attribution
methods (LRP, SHAP, LIME, Score-CAM) remain important extensions; the released
protocol is written so that each slots in without changing the pipeline. We note, however, that the
statistical and systems concerns are already addressed here: the results of
Table~\ref{tab:main} are reproduced with tight $95\%$ confidence intervals over
multiple seeds (Table~\ref{tab:seedci}), the EAQ gains are established by paired
Wilcoxon signed-rank tests (Table~\ref{tab:signif}), and per-image latency,
compression ratio and sample-level distribution moments are reported
(Tables~\ref{tab:cost} and \ref{tab:skew}). \emph{(ii)}~We study post-training compression without fine-tuning;
quantization-aware training or pruning-with-recovery may shift the knee and is
left to future work. \emph{(iii)}~The Lipschitz analysis
(Theorem~\ref{thm:lipschitz}) assumes a fixed activation pattern, which is why it
predicts the graceful regime well but only bounds---rather than models---the
collapse. \emph{(iv)}~We consider four attribution methods; LRP, SHAP and
concept-based explanations remain to be examined, though the averaging argument
of Theorem~\ref{thm:ordering} suggests conservation-respecting methods like LRP
should be comparatively robust. \emph{(v)}~Activation quantization, mixed
weight--activation schemes, and structured/hardware-aware pruning are not covered
and may interact differently with attributions.

\subsection{Theoretical Assumptions versus Physical Realities}
\label{sec:assumptions}
Table~\ref{tab:assumptions} makes the gap between our idealised assumptions and
low-bit implementation explicit, together with the mitigation each admits. The
recurring pattern is that every assumption is tight in the graceful regime and
fails in an identifiable, mechanistic way in the collapse regime---so the theory
is predictive where it claims to be and correctly flags where it does not.

\begin{table*}[!t]
\caption{Mapping of theoretical assumptions to their low-bit failure modes and
mitigations.}
\label{tab:assumptions}
\centering
\begin{tabular}{p{3.1cm} p{3.6cm} p{4.0cm} p{4.4cm}}
\hline
\textbf{Assumption} & \textbf{Manifestation} & \textbf{Failure mode ($b<3$)} &
\textbf{Mitigation / discussion}\\
\hline
Piecewise-smooth Lipschitz net (Assum.~\ref{ass:net}) &
$\sigma'(\tilde z_l)=\sigma'(z_l)$; fixed active region &
Large low-bit perturbations cross ReLU thresholds, flipping activation states and
breaking gradient continuity &
Remark~\ref{rem:reluflip} identifies ReLU sign-flipping as the explicit collapse
driver, bounding the linear-stability regime\\
High-resolution error scaling & $\mathcal D_l\!\approx\!c\sigma_l^2 2^{-2b_l}$ &
Per-layer decay exponents $\beta_l$ vary (typ.\ $3.6$--$5.3$), causing
distortion-model mismatch &
Estimate $\beta_l$ by probe quantization and use the generalised allocation
\eqref{eq:alloc-gen} (Section~\ref{sec:eaq-general})\\
Additive layer-wise distortion & $\mathcal D(\bm b)=\sum_l\omega_l2^{-2b_l}$ &
Cascading errors propagate non-linearly, giving super-additive distortion at low
precision &
Framed as a first-order Taylor model (Rem.~\ref{rem:firstorder}), exact in the
graceful regime $b\ge5$\\
Synthetic localisation ground truth & Binary shapes on low-frequency backgrounds
& Natural images add high-frequency, overlapping semantics and baseline
attribution noise &
Supplement GT masks with model-vs-model $F_\rho$/$F_{\mathrm{SSIM}}$ on
CIFAR-scale data (Section~\ref{sec:setup})\\
\hline
\end{tabular}
\end{table*}

\subsection{EAQ Ablation Study}
\label{sec:ablation}
To isolate the sources of improvement in EAQ, we conduct an ablation study on the bit allocation strategy in Table~\ref{tab:ablation}. We compare the default EAQ (which uses first-order Taylor sensitivity weights $\hat{\omega}_l$) against: (i) a \textbf{Random Allocation} that assigns random sensitivity weights to layers, and (ii) a \textbf{Magnitude Allocation} that determines layer bit-widths based on layer-wise parameter magnitudes $\norm{\bW_l}_2$. 

Random allocation leads to a drop in explanation fidelity ($0.9400$ for Grad-CAM), which is lower than uniform 5-bit quantization. Magnitude-based allocation improves upon random but still performs worse than EAQ ($0.9688$ vs $0.9715$ for Grad-CAM). This confirms that both parameter magnitude and gradient-based loss sensitivity are crucial components of the EAQ sensitivity scoring.

\begin{table}[!h]
\caption{EAQ Bit Allocation Ablation Study on CIFAR-10 (Spearman $\rho$ at 5-bit Average Budget)}
\label{tab:ablation}
\centering
\begin{tabular}{lccc}
\hline
\textbf{Method} & \textbf{Random} & \textbf{Magnitude} & \textbf{EAQ (Taylor)} \\
\hline
Saliency & 0.5191 & 0.5836 & \textbf{0.5879} \\
IG & 0.7250 & 0.7913 & \textbf{0.8018} \\
Grad-CAM & 0.9400 & 0.9688 & \textbf{0.9715} \\
\hline
\end{tabular}
\end{table}

\subsection{Limitations and Failure Cases}
\label{sec:failures}
While EAQ achieves strong results, we identify three critical failure cases and limitations:
\begin{itemize}
\item \textbf{Aggressive Compression Regimes ($b \le 2$):} When the average bit budget is extremely low, the water-filling allocation assigns $2$ bits to almost all layers. In this sub-3-bit regime, the first-order Taylor approximation breaks down because weight perturbations are too large, leading to severe activation sign-flips and dead ReLUs. In this case, EAQ fails to prevent the abrupt explanation fidelity collapse.
\item \textbf{Gradient Noise in Calibration:} Since EAQ relies on calibration gradients to estimate $\hat{\omega}_l$, out-of-distribution calibration samples or high-frequency input noise can corrupt the gradient estimates, leading to sub-optimal bit allocations.
\item \textbf{Ill-conditioned Backbones:} In networks trained without weight decay or spectral regularisation, the gradient amplification factor $\Lambda_{\mathrm{S}}$ is extremely large. This causes errors to cascade non-linearly, rendering the additive layer-wise distortion model of EAQ less accurate.
\end{itemize}

\subsection{Future Directions}
Natural extensions include: characterising fidelity under
quantization-aware training; extending EAQ to jointly allocate weight and
activation precision via a two-dimensional water-filling; deriving tighter,
distribution-specific constants $\kappa$ from the weight spectrum; and combining
EAQ with relevance-based pruning so that a single objective preserves accuracy,
sparsity and explanation fidelity simultaneously. Finally, a human-subject study
could calibrate the automated fidelity thresholds against auditors' actual
ability to detect explanation drift.

\section{Conclusion}
\label{sec:conclusion}
We presented what is, to the best of our knowledge, the first systematic
theoretical and empirical study of how neural network compression affects
post-hoc explanations. Casting quantization and
pruning as bounded weight perturbations, we derived closed-form perturbation
bounds, a Lipschitz stability theorem, a robustness ordering proving that pooled
methods (Grad-CAM, occlusion, IG) are more stable than vanilla saliency, and a
fidelity--bit-width law $\mathcal F(b)\ge1-\kappa 4^{-b}$. Experiments on a
controlled benchmark with ground-truth masks confirmed every prediction:
explanations degrade gracefully to $\sim\!5$~bits and collapse below $3$~bits,
accuracy conceals this transition, and the empirical law matches theory with
$\kappa\approx12.9$. Guided by the analysis, our Explanation-Aware Quantization
allocates precision by layer sensitivity via reverse water-filling and improves
both accuracy (up to $+6.1$\,pp) and explanation fidelity (up to $+0.062$) over
uniform quantization at a matched budget. These results argue that explanation
fidelity should join accuracy and efficiency as a first-class objective in the
design of compressed, deployable AI.

\subsection{Reproducibility and Code Availability}
\label{sec:repro}
The complete code---dataset generator, model, compression operators, all four
attribution methods, the full metric suite, and the EAQ implementation---together
with the raw per-configuration result tables and the scripts that produce every
figure in this paper are released to enable exact reproduction. All experiments
are deterministic given the fixed seeds reported in Section~\ref{sec:setup} and
run end-to-end on a single CPU core.
